% ID: SEL-D0011-BA-E02
% ESTADO: validada. S012; tres pasadas del mismo autor, sin segundo revisor.
\textbf{a) Falta de inversa.}\par\nopagebreak[4]
\[
\det A=8a-24=8(a-3).
\]
Una matriz cuadrada no tiene inversa cuando su determinante es cero: \boxed{a=3}.

\textbf{b) Ecuación para a=3.}\par\nopagebreak[4]
Aunque $A$ no sea invertible, puede serlo $A-B$. Factorizamos por la derecha, respetando el orden de los factores:
\[
XA-XB=X(A-B)=C,\qquad
A-B=\begin{pmatrix}1&2\\3&5\end{pmatrix}.
\]
Su determinante es $5-6=-1$, luego
\[
(A-B)^{-1}=\begin{pmatrix}-5&2\\3&-1\end{pmatrix}.
\]
$X$ es una fila $1\times2$. Multiplicamos a la derecha:
\[
X=C(A-B)^{-1}=(1\ \ 2)\begin{pmatrix}-5&2\\3&-1\end{pmatrix}
=\boxed{(1\ \ 0).}
\]
Comprobación: $XA-XB=(3\ \ 4)-(2\ \ 2)=(1\ \ 2)=C$. La solución es única porque $A-B$ es invertible.

\textbf{c) Identidad y potencia.}\par\nopagebreak[4]
\[
A^2=\begin{pmatrix}3&4\\6&8\end{pmatrix}^2
=\begin{pmatrix}33&44\\66&88\end{pmatrix}
=11\begin{pmatrix}3&4\\6&8\end{pmatrix}=11A.
\]
Por inducción, si $A^n=11^{n-1}A$, entonces $A^{n+1}=11^{n-1}A^2=11^nA$. Con base $n=1$, resulta
\[
\boxed{A^8=11^7A.}
\]
El exponente del escalar es $7$, no $8$; cada nueva multiplicación aporta un factor $11$.
